\documentclass[11pt]{amsart}

\usepackage[T1]{fontenc}
\usepackage{lmodern}
\usepackage{microtype}
\usepackage{amsmath,amssymb,amsthm}
\usepackage{array}
\usepackage[hidelinks]{hyperref}

\newtheorem{theorem}{Theorem}
\newtheorem{proposition}[theorem]{Proposition}
\newtheorem{lemma}[theorem]{Lemma}
\theoremstyle{definition}
\newtheorem{remark}[theorem]{Remark}

\title[Broadcast densities of the truncated square tiling]{Periodic Broadcasts of Density
\texorpdfstring{$1/10$}{1/10}, \texorpdfstring{$1/8$}{1/8} and \texorpdfstring{$5/36$}{5/36}:
Confirming and Strengthening Three Conjectured Density Bounds of Cervantes and Harris}
\date{}
\subjclass[2020]{Primary 05C69; Secondary 05C63, 94A12, 52C20}
\keywords{broadcast domination, $(t,r)$ broadcast, optimal density, truncated square tiling,
Archimedean lattice, periodic pattern}

\hypersetup{
  pdftitle={Periodic broadcasts of density 1/10, 1/8 and 5/36: confirming and strengthening
            three conjectured density bounds of Cervantes and Harris}
}

\begin{document}

\begin{abstract}
Cervantes and Harris conjectured that the optimal $(t,r)$ broadcast densities of the infinite
truncated square tiling graph $H_{\infty,\infty}$ satisfy
$\delta_{4,2}\leq 9/80$, $\delta_{4,3}\leq 1/7$ and $\delta_{4,4}\leq 1/6$, leaving all three
as conjectures short of a detailed proof; their witnesses appear there as figures. We exhibit
three explicit lattice-periodic broadcasts, each written as a congruence in the vertex
coordinates, of densities $1/10<9/80$, $1/8<1/7$ and $5/36<1/6$. So all three conjectured
inequalities hold, and all three conjectured upper bounds can be improved. A capping argument on
the published coordination sequence $1,3,5,8$ gives
$\delta_{4,2}\geq 1/13$, so $1/13\leq\delta_{4,2}(H_{\infty,\infty})\leq 1/10$; we determine
none of the three optimal densities. We note the honest baseline: since every $(3,1)$ broadcast
is a $(4,2)$ broadcast, $\delta_{4,2}\leq 1/8$ was already implied by a theorem of the same
paper, so at $r=2$ the gain over the bound the source itself implies is from $1/8$ to $1/10$.
All exhibited objects are finite and every claim about them is checkable by hand from
the tables printed here.
\end{abstract}

\maketitle

\section{The conjecture and what is proved}

Let $G$ be a graph, $t\geq 2$ and $r\geq 1$ integers. Following Blessing, Johnson, Mauretour and
Insko \cite{BJMI} in the formulation quoted by Cervantes and Harris
\cite{CervantesHarris}, a set $T\subseteq V(G)$ is a \emph{$(t,r)$ broadcast} if every
$u\in V(G)$ satisfies
\begin{equation}\label{eq:reception}
  f(u)\;:=\!\!\sum_{\substack{v\in T\\ d(u,v)<t}}\!\!\bigl(t-d(u,v)\bigr)\;\geq\;r .
\end{equation}
The distance cutoff in \eqref{eq:reception} is strict, so at $t=4$ a tower contributes
$4,3,2,1$ to the vertices at distance $0,1,2,3$ from it and nothing at distance $\geq 4$. The
\emph{density} of $T$ is $\lim_{n\to\infty}|T\cap V(H_n)|/|V(H_n)|$ along an exhausting
sequence of finite subgraphs, and $\delta_{t,r}(G)$ is the infimum of the density over all
$(t,r)$ broadcasts; Cervantes and Harris take this definition from \cite{DHR}.

Let $H_{\infty,\infty}$ denote the infinite \emph{truncated square tiling graph}, the
$3$-regular Archimedean $4.8.8$ net. Cervantes and Harris determine $\delta_{t,r}$ for
$(t,r)\in\{(2,1),(2,2),(3,1),(3,2),(3,3),(4,1)\}$ and then conjecture \cite{CervantesHarris}:

\begin{quote}
\emph{The optimal densities of $H_{\infty,\infty}$ satisfy}
\[
  \delta_{4,2}(H_{\infty,\infty})\leq \frac{9}{80},\quad
  \delta_{4,3}(H_{\infty,\infty})\leq \frac{1}{7},\quad\text{ and }\quad
  \delta_{4,4}(H_{\infty,\infty})\leq \frac{1}{6}.
\]
\end{quote}

The conjecture is a bound in each of its three coordinates, not an equality. Its authors write
that they ``provide broadcasts that achieve the bounds'' in three figures, but that ``given the
many cases needed, and short of a detailed proof, we leave the statements as conjectures''
\cite{CervantesHarris}; the accompanying problem of that paper asks for a proof that the sets
\emph{of those figures} dominate $H_{\infty,\infty}$ with the stated densities. Because
$\delta_{t,r}$ is an infimum, each inequality is an existence statement, so a single explicit
periodic broadcast settles it, and a broadcast of smaller density improves it.

\begin{theorem}\label{thm:beat}
\[
  \delta_{4,2}(H_{\infty,\infty})\leq \tfrac{1}{10}<\tfrac{9}{80},\qquad
  \delta_{4,3}(H_{\infty,\infty})\leq \tfrac{1}{8}<\tfrac{1}{7},\qquad
  \delta_{4,4}(H_{\infty,\infty})\leq \tfrac{5}{36}<\tfrac{1}{6},
\]
witnessed by the lattice-periodic broadcasts $T_5$, $T_6$, $T_7$ of \S\ref{sec:beat}. In
particular all three conjectured inequalities hold, and each of the three conjectured upper
bounds can be improved.
\end{theorem}

\begin{theorem}\label{thm:lower}
$\delta_{4,r}(H_{\infty,\infty})\geq r/C(r)$ where
$C(r)=\min(4,r)+3\min(3,r)+5\min(2,r)+8\min(1,r)$. In particular
$\delta_{4,2}\geq 2/26=1/13$, $\delta_{4,3}\geq 1/10$ and $\delta_{4,4}\geq 4/31$.
\end{theorem}

Combining Theorems \ref{thm:beat} and \ref{thm:lower},
\[
  \tfrac{1}{13}=0.076923\ldots\;\leq\;\delta_{4,2}(H_{\infty,\infty})\;\leq\;\tfrac{1}{10}=0.1,
\]
with the conjectured $9/80=0.1125$ lying strictly above the upper end. We determine none of the
three optimal densities.

We state one baseline explicitly, because it is easy to overlook.

\begin{proposition}\label{prop:baseline}
Every $(3,1)$ broadcast of a graph is a $(4,2)$ broadcast, so
$\delta_{4,2}\leq\delta_{3,1}$. Consequently $\delta_{4,2}(H_{\infty,\infty})\leq 1/8$ already
follows from the theorem of \cite{CervantesHarris} on $\delta_{3,1}$, which gives
$\delta_{3,1}(H_{\infty,\infty})\leq 1/8$.
\end{proposition}

\begin{proof}
Let $T$ be a $(3,1)$ broadcast and $u$ a vertex. Then $f_3(u)\geq 1$, so some $v\in T$ has
$3-d(u,v)\geq 1$ with $d(u,v)<3$, hence $d(u,v)\leq 2$. At $t=4$ that one tower alone
contributes $4-d(u,v)\geq 2$.
\end{proof}

So for $r=2$ the honest reading of Theorem~\ref{thm:beat} is a chain
$1/8\to 9/80\to 1/10$: the conjectured bound already strengthens the one implied by the source
paper, and $1/10$ strengthens it further. The same argument applied to the source's theorems on
$\delta_{3,2}$ and $\delta_{3,3}$ gives $\delta_{4,3}\leq\delta_{3,2}\leq 1/6$ and
$\delta_{4,4}\leq\delta_{3,3}\leq 1/4$, both weaker than the corresponding conjectured values.

\section{Coordinates, shell sizes, and the density of a periodic set}\label{sec:graph}

We use the normal form of \cite{CervantesHarris}: with the origin fixed, every vertex of
$H_{\infty,\infty}$ is a pair $(a,(x,y))$ with $a\in\mathbb{Z}_4$ and $(x,y)\in\mathbb{Z}^2$,
where $(x,y)$ indexes the squares of the underlying square tiling and $a$ numbers the four
vertices of that square clockwise from the top. We write $a@(x,y)$. Adjacency is given by six
rules: the four sides of each square,
\[
  0@(x,y)\sim 1@(x,y),\qquad 1@(x,y)\sim 2@(x,y),
\]
\[
  2@(x,y)\sim 3@(x,y),\qquad 3@(x,y)\sim 0@(x,y),
\]
and the two edge stubs that join neighbouring squares,
\[
  0@(x,y)\sim 2@(x,y+1),\qquad 1@(x,y)\sim 3@(x+1,y).
\]
The graph is $3$-regular and the translation action of $\mathbb{Z}^2$,
$a@(x,y)\mapsto a@(x+p,y+q)$, is by automorphisms.

The vertex figure of the truncated square tiling --- one square and two octagons, with no cycle
of length $3$, $5$, $6$ or $7$ through a vertex --- is a local condition, and the accompanying
program confirms it at $36$ sampled vertices of the presentation above; we do not verify that no
other Archimedean net has that vertex figure.

Only one numerical fact about $H_{\infty,\infty}$ is needed below. From each of the four vertex
types the numbers of vertices at distance $0,1,2,3,4$ are
\begin{equation}\label{eq:shells}
  1,\;3,\;5,\;8,\;11,
\end{equation}
so $|B_2(v)|=9$ and $|B_3(v)|=17$ for every $v$. These are the source's own integers: the
coordination sequence $3,5,8,11,\dots$ is quoted there from \cite{OEIS}, and the count
$|P_v(4)|=1+3+5+8=17$ appears verbatim in a figure caption of \cite{CervantesHarris}.

Finally, the device that makes every claim below finite. Let $L\subseteq\mathbb{Z}^2$ be a
sublattice of index $s$ and let $T$ be a set of vertices invariant under $L$. Then $T$ meets a
fundamental domain of $4s$ vertices in some number $k$ of them, and since the boundary of a
growing window is $O(\text{perimeter})$,
\begin{equation}\label{eq:density}
  \text{density}(T)=\frac{k}{4s}\quad\text{exactly, with no boundary correction},
\end{equation}
while $f$ is constant on each of the $4s$ classes because $L$ acts by automorphisms preserving
$T$. So verifying that an $L$-periodic $T$ is a $(t,r)$ broadcast is a finite check: evaluate
\eqref{eq:reception} once per class, reading membership of $T$ from its defining congruence.
Each of the three broadcasts below is presented that way, with its full table of $f$-values, so
no computation is left to the reader.

\section{The three witnesses}\label{sec:beat}

Throughout, $L=\langle(A,0),(B,D)\rangle$ has index $AD$ and we use the class invariants stated
with each pattern; invariance of the invariant under both generators is immediate and is checked
in each case. The labelling $T_5,T_6,T_7$ is that of the accompanying program, which also checks
patterns not needed here.

\subsection*{$T_5$: $\delta_{4,2}\leq 1/10$}
Let $L_5=\langle(5,0),(4,1)\rangle$, of index $5$, so a fundamental domain has $20$ vertices,
and put $i(x,y)=(x+y)\bmod 5$, which is $L_5$-invariant since both generators shift $x+y$ by
$5\equiv 0$. Set
\[
  \boxed{\;T_5=\bigl\{\,0@(x,y)\;:\;x+y\equiv 0\ (5)\,\bigr\}\;\cup\;
              \bigl\{\,2@(x,y)\;:\;x+y\equiv 3\ (5)\,\bigr\}.\;}
\]
This is $2$ of $20$ classes, so the density is exactly $2/20=1/10$, against the conjectured
$9/80$. The complete table of $f$, with $u=a@(i,0)$:

\smallskip
\noindent
\begin{tabular}{r|ccccc}
 & $i=0$ & $i=1$ & $i=2$ & $i=3$ & $i=4$\\
\hline
$a=0$ & $4^{\ast}$ & 2 & 4 & 2 & 3\\
$a=1$ & 4 & 3 & 4 & 3 & 2\\
$a=2$ & 2 & 4 & 2 & $4^{\ast}$ & 3\\
$a=3$ & 3 & 4 & 3 & 4 & 2
\end{tabular}
\smallskip

\noindent
The minimum is $2=r$ and the maximum is $4$, so $T_5$ is a $(4,2)$ broadcast of density $1/10$.

\subsection*{$T_6$: $\delta_{4,3}\leq 1/8$}
$L_6=\langle(4,0),(1,1)\rangle$ of index $4$ ($16$ vertices), invariant $c=(x-y)\bmod 4$,
\[
  T_6=\bigl\{\,0@(x,y):x-y\equiv 0\,\bigr\}\cup\bigl\{\,3@(x,y):x-y\equiv 2\,\bigr\}\pmod 4,
\]
which is $2$ of $16$ classes, density $1/8<1/7$. The complete table, $u=a@(c,0)$:

\smallskip
\noindent
\begin{tabular}{r|cccc}
 & $c=0$ & $c=1$ & $c=2$ & $c=3$\\
\hline
$a=0$ & $4^{\ast}$ & 4 & 4 & 4\\
$a=1$ & 3 & 4 & 3 & 5\\
$a=2$ & 3 & 5 & 3 & 4\\
$a=3$ & 4 & 4 & $4^{\ast}$ & 4
\end{tabular}
\smallskip

\noindent
The minimum over the sixteen classes is $3=r$, so $T_6$ is a
$(4,3)$ broadcast of density $1/8$.

\subsection*{$T_7$: $\delta_{4,4}\leq 5/36$}
$L_7=\langle(6,0),(0,3)\rangle$ of index $18$, so a fundamental domain has $72$ vertices, with
class invariant $(x\bmod 6,\,y\bmod 3)$. Take the ten classes
\begin{align*}
  T_7=\bigl\{\;&0@(0,0),\;0@(2,0),\;0@(4,1),\;1@(0,2),\;1@(3,0),\\
               &2@(1,1),\;2@(3,2),\;2@(5,2),\;3@(2,2),\;3@(5,0)\;\bigr\},
\end{align*}
each entry read as ``all $a@(x,y)$ with $x\equiv\cdot\pmod 6$ and $y\equiv\cdot\pmod 3$''. This
is $10$ of $72$ classes, density $10/72=5/36=0.13888\ldots<1/6$. All seventy-two values of $f$,
with the six blocks of columns indexed by $x\bmod 6$ and the three columns of each block by
$y\bmod 3$:

\smallskip
\noindent
\footnotesize
\setlength{\tabcolsep}{3pt}
\begin{tabular}{r|ccc|ccc|ccc|ccc|ccc|ccc}
 & \multicolumn{3}{c|}{$x\equiv 0$} & \multicolumn{3}{c|}{$x\equiv 1$}
 & \multicolumn{3}{c|}{$x\equiv 2$} & \multicolumn{3}{c|}{$x\equiv 3$}
 & \multicolumn{3}{c|}{$x\equiv 4$} & \multicolumn{3}{c}{$x\equiv 5$}\\
$y\equiv$ & 0 & 1 & 2 & 0 & 1 & 2 & 0 & 1 & 2 & 0 & 1 & 2 & 0 & 1 & 2 & 0 & 1 & 2\\
\hline
$a=0$ & $4^{\ast}$ & 4 & 5 & 5 & 4 & 4 & $4^{\ast}$ & 4 & 5 & 4 & 4 & 4 & 5 & $4^{\ast}$ & 5 & 4 & 4 & 4\\
$a=1$ & 5 & 5 & $4^{\ast}$ & 4 & 4 & 4 & 5 & 4 & 4 & $4^{\ast}$ & 5 & 5 & 4 & 4 & 4 & 4 & 4 & 5\\
$a=2$ & 4 & 4 & 4 & 5 & $4^{\ast}$ & 5 & 4 & 4 & 4 & 5 & 4 & $4^{\ast}$ & 4 & 4 & 5 & 5 & 4 & $4^{\ast}$\\
$a=3$ & 5 & 4 & 4 & 4 & 4 & 4 & 5 & 5 & $4^{\ast}$ & 4 & 4 & 5 & 4 & 4 & 4 & $4^{\ast}$ & 5 & 5
\end{tabular}
\normalsize
\setlength{\tabcolsep}{6pt}
\smallskip

\noindent
Every entry is at least $4=r$ and the maximum is $5$, so $T_7$ is a $(4,4)$ broadcast. This
proves Theorem~\ref{thm:beat}.

\section{The lower bound}\label{sec:lower}

\begin{lemma}[capping]\label{lem:cap}
Let $T$ be a $(4,r)$ broadcast and set
\[
  g(u)=\sum_{v\in T,\;d(u,v)\leq 3}\min\bigl(4-d(u,v),\,r\bigr).
\]
Then $g(u)\geq r$ for all $u$.
\end{lemma}

\begin{proof}
If some single $v\in T$ has $4-d(u,v)\geq r$ then its own term is
$\min(4-d(u,v),r)=r$ and we are done. Otherwise $4-d(u,v)<r$ for every tower in the ball, so
every term is unchanged by the cap and $g(u)=f(u)\geq r$.
\end{proof}

\begin{proof}[Proof of Theorem~\ref{thm:lower}]
By \eqref{eq:shells} each tower $v$ contributes $\min(4-d,r)$ to each of the $1,3,5,8$ vertices
at distance $d=0,1,2,3$ from it, so it contributes exactly
$C(r)=\min(4,r)+3\min(3,r)+5\min(2,r)+8\min(1,r)$ to $\sum_u g(u)$. Summing
Lemma~\ref{lem:cap} over the vertices of a window $H_n$ and dividing by $|V(H_n)|$ gives
$r\leq\bigl(|T\cap V(H_n)|/|V(H_n)|\bigr)C(r)+O(|\partial H_n|/|V(H_n)|)$, and the error term
vanishes. Hence $\text{density}(T)\geq r/C(r)$ for every $(4,r)$ broadcast. Finally
$C(2)=2+6+10+8=26$, $C(3)=3+9+10+8=30$ and $C(4)=4+9+10+8=31$.
\end{proof}

\section{What is not settled}\label{sec:open}

\begin{itemize}
\item \textbf{The optimal densities.} The gaps
$\delta_{4,2}\in[1/13,1/10]$, $\delta_{4,3}\in[1/10,1/8]$ and
$\delta_{4,4}\in[4/31,5/36]$ are open. Nothing above shows either end is attained; in
particular nothing above asserts $\delta_{4,2}=1/10$, and since $\delta_{4,r}$ is an infimum an
optimal broadcast need not be periodic at all.
\item \textbf{The literal form of the companion problem of \cite{CervantesHarris}.} That
problem asks for a proof that the sets \emph{of the paper's own three figures} dominate
$H_{\infty,\infty}$ with the stated densities. Those sets appear in the source only as raster
images and we did not transcribe them; the patterns above are independent objects and need not
coincide with them. We therefore assert nothing whatever about the authors' figures --- in
particular no claim about which reception mechanism they use. The inequalities of the conjecture
follow from Theorem~\ref{thm:beat}, read with the reception rule \eqref{eq:reception} and the
density of \S\ref{sec:graph}; the literal reading of that problem is untouched. The further
companion question of \cite{CervantesHarris} concerns the six pairs $(t,r)$ for which
$\delta_{t,r}$ is proved there, not $(4,2)$, $(4,3)$ or $(4,4)$, and is not addressed here.
\item \textbf{Adjacent work.} The efficient- and perfect-domination results for Archimedean
lattices of Zhao, Wierman and Marge \cite{ZWM2018,ZWM2022} concern domination ratios rather than
$(t,r)$ broadcast densities, and give none of the bounds above. The $(t,r)$ broadcast literature
closest in parameters works on the infinite square grid and its higher-dimensional analogues
rather than on $H_{\infty,\infty}$: Blessing, Johnson, Mauretour and Insko \cite{BJMI}, Drews,
Harris and Randolph \cite{DHR} together with their computational companion \cite{DHRupper},
Herrman and van Hintum \cite{HvH}, and Shlomi \cite{Shlomi}. Those papers bound broadcast
densities of grids, and none of them bounds $\delta_{4,2}$, $\delta_{4,3}$ or $\delta_{4,4}$ of
the truncated square tiling. In the published version of \cite{CervantesHarris} the three
inequalities above are still stated as a conjecture, with the same three values, and its problem
list still asks for a proof.
\end{itemize}

\section*{Verification}

The exhibited objects and the hand arithmetic are re-derived from the congruences printed here by
the accompanying program \texttt{verify.py}, which rebuilds the graph from the six adjacency
rules of \S\ref{sec:graph}, performs a fresh breadth-first search in the infinite graph at every
class representative, and uses exact rational arithmetic throughout.

The statement the program bears on is Theorem~\ref{thm:beat}. The checks that carry it are the
presentation of the graph and the shell sizes \eqref{eq:shells}, and, for each of
$T_5,T_6,T_7$, the lattice invariance, the class count, the density and all $20+16+72$ typeset
$f$-values, together with the comparisons $1/10<9/80$, $1/8<1/7$ and $5/36<1/6$; in the
transcript these are the lines prefixed \texttt{graph}, \texttt{shells}, \texttt{ball-sizes},
\texttt{W5}, \texttt{W6} and \texttt{W7}. Theorem~\ref{thm:lower} is \emph{not} verified there:
the program recomputes only the integers $C(2)=26$, $C(3)=30$, $C(4)=31$ and the fractions
$r/C(r)$, and its own lines record that the capping step is not checked. So
Theorem~\ref{thm:lower}, and with it the left-hand end of the sandwich
$1/13\leq\delta_{4,2}\leq 1/10$, rests on Lemma~\ref{lem:cap} and the hand proof of
\S\ref{sec:lower} alone. Nor does the program establish Proposition~\ref{prop:baseline}: it runs
the source's own $(3,1)$ broadcast of density $1/8$ through the $(3,1)$ and then the $(4,2)$
predicate, which is that one witness and not the general implication; the implication is proved
above by hand.

The program's closing line reads \texttt{VERDICT: ALL 99 CHECKS PASS}, but those $99$ include
checks on further patterns and comparisons this paper does not use, so that line is not by itself
a certificate of the statements above: only the groups named in the previous paragraph are. The
program does not
read the source paper's raster figures, which were never transcribed, so it makes no claim about
them; it makes no claim about optimality either.

The tables are also short enough to check by hand, and several entries were.

\begin{thebibliography}{9}

\bibitem{BJMI}
D.~Blessing, K.~Johnson, C.~Mauretour and E.~Insko,
\emph{On $(t,r)$ broadcast domination numbers of grids},
Discrete Appl. Math. \textbf{187} (2015), 19--40.
\href{https://doi.org/10.1016/j.dam.2015.02.005}{doi:10.1016/j.dam.2015.02.005}

\bibitem{CervantesHarris}
J.~Cervantes and P.~E.~Harris,
\emph{$(t,r)$ broadcast domination numbers and densities of the truncated square tiling graph},
Res. Math. Sci. \textbf{13} (2026), no.~2.
\href{https://doi.org/10.1007/s40687-026-00623-0}{doi:10.1007/s40687-026-00623-0};
preprint \href{https://arxiv.org/abs/2408.13331}{arXiv:2408.13331v1}, 2024.
Every quotation above was read from the arXiv \LaTeX{} e-print, in which the statements cited
here occur at the following source lines: the density conjecture at 1292--1298, the companion
problem on its three figures at 1417--1419, the further companion question at 1422, and the
theorems on $\delta_{3,1}$, $\delta_{3,2}$ and $\delta_{3,3}$ at 666, 813 and 1077
respectively. In the published version the density conjecture is Conjecture 4.10, the problem
asking for a proof for its own three figures is Problem 5, and the bounds
$\delta_{3,1}\leq 1/8$ and $\delta_{3,2}\leq 1/6$ are Theorems 4.5 and 4.6.

\bibitem{DHR}
B.~F.~Drews, P.~E.~Harris and T.~W.~Randolph,
\emph{Optimal $(t,r)$ broadcasts on the infinite grid},
Discrete Appl. Math. \textbf{255} (2019), 183--197.
\href{https://doi.org/10.1016/j.dam.2018.08.009}{doi:10.1016/j.dam.2018.08.009}

\bibitem{DHRupper}
B.~F.~Drews, P.~E.~Harris and T.~W.~Randolph,
\emph{Computing upper bounds for optimal density of $(t,r)$ broadcasts on the infinite grid},
preprint \href{https://arxiv.org/abs/1712.00150}{arXiv:1712.00150}, 2017.

\bibitem{HvH}
R.~A.~Herrman and P.~van Hintum,
\emph{$(t,r)$ broadcast domination in the infinite grid},
preprint \href{https://arxiv.org/abs/1912.11560}{arXiv:1912.11560}, 2019.

\bibitem{Shlomi}
E.~Shlomi,
\emph{Bounds on $(t,r)$ broadcast domination of $n$-dimensional grids},
Discrete Math. Theor. Comput. Sci. \textbf{24} (2023), no.~1;
preprint \href{https://arxiv.org/abs/1908.07586}{arXiv:1908.07586}, 2019.

\bibitem{OEIS}
OEIS Foundation Inc.,
\emph{Entry A008576: coordination sequence for the $4.8.8$ Archimedean net},
The On-Line Encyclopedia of Integer Sequences,
\href{https://oeis.org/A008576}{oeis.org/A008576}

\bibitem{ZWM2018}
Y.~Zhao, J.~C.~Wierman and T.~G.~Marge,
\emph{Efficient and non-efficient domination of Archimedean lattices},
Congr. Numer. \textbf{231} (2018), 95--108.

\bibitem{ZWM2022}
Y.~Zhao, J.~C.~Wierman and T.~G.~Marge,
\emph{Perfect domination ratios of Archimedean lattices},
Electron. J. Combin. \textbf{29} (2022), no.~3, Paper No.~3.60.
\href{https://doi.org/10.37236/10210}{doi:10.37236/10210};
see also Congr. Numer. \textbf{231} (2018), 39--61.

\end{thebibliography}

\end{document}
